\documentclass[a4paper,12pt]{article}
\usepackage[T2A]{fontenc}
\usepackage[utf8]{inputenc}
\usepackage[russian]{babel}
\usepackage{amsmath,amssymb,mathtools}
\usepackage{icomma}
\usepackage[a4paper,left=21mm,right=21mm,top=22mm,bottom=21mm,headheight=26pt]{geometry}
\usepackage{microtype}
\usepackage{tikz}
\usetikzlibrary{arrows.meta,positioning,shapes.geometric,calc}
\usepackage{tabularx,booktabs,array}
\usepackage{enumitem}
\usepackage{parskip}
\usepackage{fancyhdr}
\usepackage{setspace}
\usepackage{needspace}
\usepackage{xcolor}
\usepackage[hidelinks]{hyperref}
\definecolor{accent}{HTML}{008080}
\definecolor{darkaccent}{HTML}{00585A}
\definecolor{textgray}{HTML}{333333}
\definecolor{muted}{HTML}{68777A}
\definecolor{softline}{HTML}{C5D6D4}
\color{textgray}
\onehalfspacing
\setlength{\parindent}{0pt}
\setlength{\parskip}{5pt}
\setlist[enumerate]{leftmargin=*,itemsep=4pt,topsep=3pt}
\setlist[itemize]{leftmargin=*,itemsep=3pt,topsep=2pt}
\newcommand{\key}[1]{\textcolor{darkaccent}{\textbf{#1}}}
\newcommand{\checkitem}[1]{\noindent\textcolor{accent}{\(\square\)}\hspace{0.3em}#1\par}
\newcommand{\ex}[1]{\par\medskip\textcolor{darkaccent}{\textbf{Пример. #1}}\par}
\newcommand{\topic}[1]{\section{#1}}
\setcounter{secnumdepth}{1}
\renewcommand{\thesection}{\arabic{section}}
\usepackage{titlesec}
\titleformat{\section}{\normalfont\large\color{darkaccent}}{\thesection.}{0.55em}{}
\titlespacing*{\section}{0pt}{15pt}{9pt}
\titleformat{\subsection}{\normalfont\normalsize\color{darkaccent}}{}{0pt}{}
\titlespacing*{\subsection}{0pt}{9pt}{5pt}
\pagestyle{fancy}
\fancyhf{}
\fancyhead[L]{\small\color{muted}Остатки от деления}
\fancyhead[R]{\small\color{muted}Ксения Клыкова}
\fancyfoot[C]{\small\thepage}
\renewcommand{\headrulewidth}{0.35pt}
\renewcommand{\headrule}{\hbox to\headwidth{\color{accent}\leaders\hrule height \headrulewidth\hfill}}
\hypersetup{pdftitle={Чек-лист. Остатки от деления и олимпиадные задачи},pdfauthor={Лёвин Артём Александрович},pdfsubject={Математика. Деление с остатком}}
\tikzset{
  endstep/.style={ellipse,draw=accent,line width=.8pt,minimum width=37mm,minimum height=10mm,align=center,inner sep=3pt},
  actstep/.style={rectangle,rounded corners=2pt,draw=softline,line width=.85pt,text width=101mm,minimum height=13mm,align=center,inner sep=5pt},
  teststep/.style={diamond,aspect=2.8,draw=accent,line width=.9pt,align=center,text width=56mm,inner sep=2pt},
  arr/.style={-{Stealth[length=2.6mm,width=1.8mm]},line width=.9pt,color=darkaccent}
}
\begin{document}
\thispagestyle{empty}
\begin{center}
\vspace*{25mm}
{\large\color{darkaccent} МАТЕМАТИКА}\par
\vspace{21mm}
{\LARGE Чек-лист}\par
\vspace{13mm}
{\Large\color{accent} Остатки от деления\\[3pt] и олимпиадные задачи}\par
\vspace{20mm}
{\normalsize Олимпиадная подготовка по математике}\par
\vfill
{\large Ксения Клыкова}\par
\vspace{14mm}
{20 ноября 2026 г.}\par
\vspace{19mm}
{\small Лёвин Артём Александрович\\[-1pt]эксклюзивно для Ксении Клыковой}\par
\vspace{19mm}
\begin{tikzpicture}[x=1cm,y=1cm]
 \draw[color=accent,line width=0.8pt] (0,0) .. controls (1.2,0.55) and (2.3,-0.42) .. (3.7,0.05)
 .. controls (5.0,0.56) and (6.2,-0.44) .. (7.7,0.05);
\end{tikzpicture}\par
\end{center}
\clearpage

\topic{Что означает остаток от деления}
\key{Основная формула:} для любого целого числа $n$ и положительного целого $m$ существуют единственные целые $k$ и $r$, для которых
\[
\boxed{n=mk+r,\qquad k\in\mathbb Z,\qquad 0\le r<m.}
\]
Здесь $m$ --- делитель (\emph{модуль}), $k$ --- неполное частное, $r$ --- остаток. Остаток может быть \textbf{нулевым}, но не бывает отрицательным и всегда меньше делителя.

\ex{Одно и то же основание}
\[
12=2\cdot5+2,\quad 17=3\cdot5+2,\quad 22=4\cdot5+2.
\]
Во всех трёх случаях при делении на $5$ остаток равен $2$. Изменяется только целое число $k$.

\checkitem{Умею выделять делитель, неполное частное и остаток.}
\checkitem{Проверяю два обязательных условия: $m>0$ и $0\le r<m$.}

\topic{Числа с одинаковым остатком}
Все целые числа с остатком $r$ при делении на $m$ имеют вид
\[
\boxed{n=mk+r,\qquad k\in\mathbb Z.}
\]
Два таких числа отличаются на число, кратное $m$. Поэтому следующий элемент последовательности получают прибавлением $m$:
\[
\underbrace{2}_{k=0}\ \xrightarrow{\,+5\,}\ \underbrace{7}_{k=1}\ \xrightarrow{\,+5\,}\ \underbrace{12}_{k=2}\ \xrightarrow{\,+5\,}\ \underbrace{17}_{k=3}.
\]
Дополнительная запись, полезная в олимпиадах: $n\equiv r\pmod m$ означает, что $n$ и $r$ дают одинаковый остаток при делении на $m$.

\checkitem{Записываю все подходящие числа с помощью одного целого параметра $k$.}
\checkitem{Понимаю, почему разность двух подходящих чисел кратна модулю.}

\topic{Поиск чисел на заданном отрезке}
Если $A\le n\le B$, а $n=mk+r$, подставляем формулу числа в двойное неравенство. Затем выбираем \textbf{все целые} значения $k$ между границами и возвращаемся к $n=mk+r$:
\[
A\le mk+r\le B
\quad\Longleftrightarrow\quad
\boxed{\frac{A-r}{m}\le k\le\frac{B-r}{m}}\qquad(m>0).
\]

\ex{Найдите все числа от $30$ до $60$, дающие остаток $5$ при делении на $8$}
\[
30\le8k+5\le60\quad\Longrightarrow\quad
\frac{25}{8}\le k\le\frac{55}{8}.
\]
Подходят только $k=4,5,6$. Возвращаемся к исходному числу:
\[
\boxed{n\in\{37,45,53\}.}
\]
Проверка: $37=4\cdot8+5$, $45=5\cdot8+5$, $53=6\cdot8+5$. \textbf{Всего три числа.}

\subsection*{Как быстро подсчитать количество}
Если первое и последнее подходящие числа --- $n_{\min}$ и $n_{\max}$, то
\[
\boxed{N=\frac{n_{\max}-n_{\min}}{m}+1.}
\]
В примере $N=(53-37)/8+1=3$. Единица учитывает первый элемент. Если известны только границы для $k$, можно использовать
\[
N=\max\!\left(0,\left\lfloor\frac{B-r}{m}\right\rfloor-
\left\lceil\frac{A-r}{m}\right\rceil+1\right).
\]
Здесь $\lfloor x\rfloor$ --- округление вниз до целого, $\lceil x\rceil$ --- вверх. Для пустого интервала целых $k$ ответ равен нулю.

\ex{Числа с остатком $3$ по модулю $7$ на $[-100;100]$}
\[
-100\le 7k+3\le100
\quad\Longrightarrow\quad
-\frac{103}{7}\le k\le\frac{97}{7}.
\]
Получаем $k=-14,-13,\ldots,13$, то есть
\[
N=13-(-14)+1=\boxed{28}.
\]
Крайние числа: $-95$ и $94$. Поэтому также $N=(94-(-95))/7+1=28$.

\checkitem{При выборе $k$ не округляю границы произвольно.}
\checkitem{Различаю просьбы «найдите все числа» и «найдите их количество».}
\checkitem{Не забываю прибавить $1$ в формуле количества.}

\Needspace{9\baselineskip}
\topic{Отрицательные числа}
Для отрицательных чисел используется \textbf{тот же} принцип: $n=mk+r$ и $0\le r<m$.
\[
\boxed{-251=(-21)\cdot12+1}\quad\Rightarrow\quad r=1.
\]
Запись $-251=(-20)\cdot12-11$ верна как равенство, но $-11$ \textbf{не является допустимым остатком}. При делении $-146$ на $11$:
\[
-146=(-14)\cdot11+8,\qquad \boxed{r=8}.
\]
\checkitem{Для отрицательных чисел тоже получаю остаток из диапазона $0,1,\ldots,m-1$.}

\topic{Остатки и периодические процессы}
В задачах с повторяющимся циклом \textbf{модуль равен длине периода}, а остаток показывает смещение внутри очередного цикла.

\ex{Сейчас на 12-часовом циферблате 9 часов. Что будет через 100 часов?}
\[
100=8\cdot12+4,\qquad 9+4=13.
\]
После 12 часов стрелка проходит полный круг, поэтому отметка $13$ соответствует отметке $\boxed{1}$ на 12-часовом циферблате.

\ex{Дни недели}
Неделя повторяется через $7$ дней. Например, через $45$ дней после понедельника наступит четверг, поскольку $45=6\cdot7+3$.

\checkitem{Нахожу период и использую остаток как смещение.}
\checkitem{Привожу окончательную отметку к диапазону одного цикла.}

\topic{Как изменить модуль, зная остаток}
Пусть $n=12k+7$. Чтобы найти остатки от деления $n$ на $3$, $4$ и $6$, замечаем: число $12k$ делится на каждое из этих чисел. Поэтому остаток зависит только от $7$:
\[
\begin{array}{c|ccc}
\text{Делитель} & 3 & 4 & 6\\ \midrule
\text{Остаток} & 1 & 3 & 1
\end{array}
\]
\key{Правило:} если $n=mk+r$ и $d$ делит $m$, то остаток $n$ при делении на $d$ совпадает с остатком $r$ при делении на $d$.

\Needspace{5\baselineskip}
\textit{Важно:} если $d$ не делит $m$, остаток может зависеть от $k$.

\checkitem{Проверяю, делится ли коэффициент при $k$ на новый делитель.}
\checkitem{Не отбрасываю слагаемое $mk$, пока не обосновал его делимость.}

\topic{Два условия на остатки}
\textbf{Задача:} найдите наименьшее положительное целое число $n$, которое даёт остаток $4$ при делении на $6$ и остаток $2$ при делении на $5$.

Записываем оба условия с \textbf{разными} целыми параметрами:
\[
n=6k+4,\qquad n=5t+2,\qquad k,t\in\mathbb Z.
\]
Простой путь --- перебрать числа из первой последовательности, проверяя второе условие:
\[
\begin{array}{c|rrrr}
n & 4 & 10 & 16 & \mathbf{22}\\ \midrule
\text{Остаток при делении на }5 &4&0&1&\mathbf{2}
\end{array}
\]
Искомое наименьшее положительное число равно $\boxed{22}$. Поскольку $6$ и $5$ взаимно просты, все решения отличаются на $6\cdot5=30$:
\[
\boxed{n=22+30s,\qquad s\in\mathbb Z.}
\]
\textbf{Почему $30$?} При прибавлении $30$ оба остатка сохраняются. Любая разность двух решений делится и на $6$, и на $5$, следовательно, делится на $30$.

\topic{Как перейти от двух условий к целочисленной модели}
Равенство $6k+4=5t+2$ приводит к
\[
6k=5t-2,\qquad k=\frac{5t-2}{6}.
\]
Для каждого целого $t$ выражение $(5t-2)/6$ также должно быть целым. Это пример уравнения в целых числах (диофантова уравнения).

Для подсчёта решений на отрезке $[A;B]$ составляем $A\le 22+30s\le B$.
Например, на $[-100;100]$:
\[
-\frac{122}{30}\le s\le\frac{78}{30}
\quad\Longrightarrow\quad s=-4,-3,-2,-1,0,1,2.
\]
\textbf{Ответ:} $7$ решений: $-98,-68,-38,-8,22,52,82$.

\Needspace{17\baselineskip}
\topic{Типичные ошибки и их исправление}
\begin{itemize}
\item \textbf{Остаток отрицателен или равен модулю.} Требуется $0\le r<m$.
\item \textbf{Найден параметр $k$, но не число $n$.} Вернитесь к $n=mk+r$.
\item \textbf{Потеря крайнего целого $k$.} Проверяйте обе границы неравенства, включая равенство.
\item \textbf{Забыто $+1$ при подсчёте.} Между двумя крайними элементами на один промежуток меньше, чем элементов.
\item \textbf{Цикл не завершён.} На циферблате после $12$ снова идёт $1$.
\item \textbf{Переменные при двух модулях отождествлены.} В формулах $6k+4$ и $5t+2$ частные разные.
\item \textbf{Не проверена кратность нового делителя старому.} Сначала убедитесь, что целое слагаемое действительно даёт остаток $0$.
\end{itemize}

\subsection*{Итоговая памятка}
\renewcommand{\arraystretch}{1.30}
\begin{tabularx}{\linewidth}{@{}p{.38\linewidth}X@{}}
\toprule
\textcolor{darkaccent}{\textbf{Сигнал в задаче}} & \textcolor{darkaccent}{\textbf{Ваше действие}}\\
\midrule
Остаток $r$ по модулю $m$ & Запишите $n=mk+r$, $k\in\mathbb Z$.\\
Ограниченный отрезок & Подставьте формулу в двойное неравенство.\\
Нужно количество & Считайте целые $k$ или примените $(n_{\max}-n_{\min})/m+1$.\\
Отрицательное число & Добейтесь $0\le r<m$.\\
Цикл или часы & Возьмите период как модуль.\\
Два остаточных условия & Найдите первое решение, затем общий шаг.\\
\bottomrule
\end{tabularx}

\clearpage
\thispagestyle{fancy}
\begin{center}
{\Large\color{darkaccent} Алгоритм: числа с заданным остатком на отрезке}
\end{center}
\vspace{10mm}
\begin{center}
\begin{tikzpicture}[node distance=10mm]
 \node[endstep] (start) {Начало};
 \node[actstep,below=of start] (model) {Запишите условие через целое число:\\$n=mk+r$, где $0\le r<m$};
 \node[actstep,below=of model] (bound) {Учитывайте границы отрезка:\\$A\le mk+r\le B$};
 \node[actstep,below=of bound] (solve) {Выразите целый параметр:\\$\dfrac{A-r}{m}\le k\le\dfrac{B-r}{m}$};
 \node[teststep,below=13mm of solve] (choice) {Есть целые $k$ между границами?};
 \node[actstep,below=14mm of choice] (yes) {Подставьте все целые $k$ в $n=mk+r$;\\если нужно количество --- посчитайте их};
 \node[endstep,below=of yes] (finish) {Ответ};
 \node[actstep,right=12mm of choice,text width=29mm,minimum height=14mm] (no) {Подходящих чисел нет};
 \node[endstep,below=of no,minimum width=31mm] (none) {$N=0$};
 \draw[arr] (start) -- (model);
 \draw[arr] (model) -- (bound);
 \draw[arr] (bound) -- (solve);
 \draw[arr] (solve) -- (choice);
 \draw[arr] (choice) -- node[right,near start,fill=white,inner sep=1pt]{\small да} (yes);
 \draw[arr] (yes) -- (finish);
 \draw[arr] (choice.east) -- node[above,fill=white,inner sep=1pt]{\small нет} (no.west);
 \draw[arr] (no) -- (none);
\end{tikzpicture}
\end{center}
\clearpage

\begin{center}
{\Large\color{darkaccent} Тренировочный блок}\par
\vspace{1mm}
{\small\color{muted}Путь А --- 10 задач с постепенным повышением сложности}
\end{center}
\vspace{4mm}
\textcolor{darkaccent}{\textbf{Базовый уровень}}
\begin{enumerate}
\item При делении числа $58$ на $9$ найдите неполное частное и остаток.
\item Выпишите все неотрицательные числа, меньшие $40$, которые дают остаток $3$ при делении на $7$.
\item Найдите все целые числа на отрезке $[25;80]$, дающие остаток $4$ при делении на $9$.
\end{enumerate}
\medskip
\textcolor{darkaccent}{\textbf{Средний уровень}}
\begin{enumerate}[resume]
\item Сколько целых чисел на отрезке $[-60;75]$ дают остаток $2$ при делении на $8$?
\item Сейчас на 12-часовом циферблате 10 часов. На какой отметке будет часовая стрелка через $137$ часов?
\item Сегодня понедельник. Какой день недели наступит через $45$ дней?
\item Найдите неполное частное и остаток от деления числа $-146$ на $11$ в записи $-146=11k+r$, где $0\le r<11$.
\end{enumerate}
\medskip
\textcolor{darkaccent}{\textbf{Повышенный уровень}}
\begin{enumerate}[resume]
\item Целое число имеет вид $n=18k+11$, $k\in\mathbb Z$. Найдите его остатки при делении на $2$, $3$, $6$ и $9$.
\item Найдите наименьшее положительное целое число, дающее остаток $3$ при делении на $7$ и остаток $4$ при делении на $5$.
\item Сколько целых чисел на отрезке $[-150;150]$ дают одновременно остаток $4$ при делении на $6$ и остаток $2$ при делении на $5$?
\end{enumerate}
\vfill
\textcolor{muted}{\small При решении задач на отрезке сначала найдите диапазон целых значений параметра, затем возвращайтесь к исходным числам.}
\clearpage

\begin{center}
{\Large\color{darkaccent} Ответы к тренировочному блоку}
\end{center}
\vspace{8mm}
\renewcommand{\arraystretch}{1.65}
\begin{tabularx}{\linewidth}{@{}>{\centering\arraybackslash}p{14mm}X@{}}
\toprule
\textcolor{darkaccent}{\textbf{№}} & \textcolor{darkaccent}{\textbf{Ответ}}\\
\midrule
1 & $k=6$, остаток $4$.\\
2 & $3,10,17,24,31,38$.\\
3 & $31,40,49,58,67,76$.\\
4 & $17$ чисел.\\
5 & $3$ часа (отметка $3$).\\
6 & Четверг.\\
7 & $k=-14$, остаток $8$.\\
8 & По порядку: $1,2,5,2$.\\
9 & $24$.\\
10 & $10$ чисел.\\
\bottomrule
\end{tabularx}
\vspace{20mm}
\begin{center}
{\small\color{muted}Главная идея: фиксированный остаток задаёт последовательность с постоянным шагом. Ограничение отрезком превращает задачу в поиск целых значений параметра.}
\end{center}
\vfill
\begin{center}
{\small Лёвин Артём Александрович\quad\textcolor{accent}{\textbullet}\quad эксклюзивно для Ксении Клыковой}
\end{center}
\end{document}
